\documentclass[11pt,a4paper]{article}

% Packages
\usepackage[utf8]{inputenc}
\usepackage[T1]{fontenc}
\usepackage{amsmath,amssymb,amsthm}
\usepackage{mathtools}
\usepackage{graphicx}
\usepackage{hyperref}
\usepackage{cleveref}
\usepackage{booktabs}
\usepackage{array}
\usepackage{xcolor}
\usepackage{microtype}
\usepackage[numbers]{natbib}

% Page layout
\usepackage[margin=1in]{geometry}

% Theorem environments
\theoremstyle{plain}
\newtheorem{theorem}{Theorem}[section]
\newtheorem{lemma}[theorem]{Lemma}
\newtheorem{corollary}[theorem]{Corollary}
\newtheorem{proposition}[theorem]{Proposition}
\newtheorem{conjecture}[theorem]{Conjecture}

\theoremstyle{definition}
\newtheorem{definition}[theorem]{Definition}
\newtheorem{remark}[theorem]{Remark}

% Notation shortcuts
\newcommand{\F}{\mathbb{F}}
\newcommand{\LL}{\mathcal{L}}
\newcommand{\E}{\mathbb{E}}
\newcommand{\Prb}{\Pr}
\DeclareMathOperator{\Lagr}{Lagr}
\DeclareMathOperator{\SDA}{SDA}
\DeclareMathOperator{\VSTAT}{VSTAT}
\DeclareMathOperator{\Adv}{Adv}
\newcommand{\etal}{\textit{et al.}}
\newcommand{\R}{\mathbb{R}}
\newcommand{\LSN}{\mathsf{LSN}}

\title{The Lagrangian Subspace Noise Problem: \\ A New Framework for Post-Quantum Cryptography}
\author{Kwanghoo Choo\thanks{ORCID: \href{https://orcid.org/0009-0005-5682-8098}{0009-0005-5682-8098}. Email: \texttt{gattochoo@gmail.com}}}
\date{June 2026}

\begin{document}

\maketitle

\begin{abstract}
We introduce the Lagrangian Subspace Noise (LSN) problem --- a variant of the Learning Parity with Noise (LPN) problem where the secret is a Lagrangian subspace of a symplectic vector space rather than a linear function. The secret space, the Lagrangian Grassmannian $\Lagr(2n,\F_2)$, has cardinality $2^{n(n+1)/2+O(1)}$, exponentially larger than LPN's $2^n$ yet tightly constrained by the symplectic form. We prove an unconditional exponential Statistical Query (SQ) lower bound of $\Omega(2^n)$ queries via an explicit symplectic spread, and a conditional $2^{2n-O(1)}$ bound under a precisely-stated extremality conjecture, together with an exact pairwise-correlation formula and no asymptotic error term. We establish that linear and low-degree polynomial feature-map reductions are information-theoretically impossible, that adaptive linear SQ algorithms are completely blocked by direct structural analysis, and we map the exact standard-model gap that remains for general non-linear reductions. We construct two cryptographic primitives from LSN: a succinct non-interactive argument (LSN-SNARK) with $O(n^2)$ R1CS constraints ($\approx 4{,}225$ for $n=65$, compared to millions for ZK-lattice proofs), and a key-encapsulation mechanism (LSN-KEM) with concatenated polar-code reconciliation achieving 80-bit security with public keys of 1.78~KB and ciphertexts of 288~B. All security claims are explicitly classified as theorem, evidence, or conjecture, and the remaining open problems are stated precisely.
\end{abstract}

\paragraph{\textbf{Keywords.}} Post-quantum cryptography, Learning Parity with Noise, statistical query model, symplectic geometry, polar codes, zero-knowledge proofs.

\section*{Acknowledgements}
During the preparation of this work, the author used large language models (Claude, Kimi, and Codex) to assist with literature review, mathematical exploration, code prototyping, and manuscript drafting. After using these tools, the author reviewed and edited the content as needed and takes full responsibility for the accuracy, integrity, and final content of the work.

\paragraph{\textbf{Data and code availability.}} All Python prototypes, test vectors, and the \LaTeX{} source are available at \url{https://github.com/Gattochoo/lsn-pq}.

\paragraph{\textbf{Conflict of interest.}} The author declares no competing interests.

\section{Introduction}
\label{sec:intro}

\subsection{Background}
\label{subsec:background}

The Learning Parity with Noise (LPN) problem, introduced by Blum \etal\ \cite{BFKL94}, asks to recover a secret linear function $s \in \F_2^k$ given noisy inner products $(a_i, \langle a_i, s \rangle \oplus e_i)$. Despite decades of cryptanalytic effort, no polynomial-time algorithm is known for LPN with constant noise rate. The best classical attacks run in sub-exponential time $2^{O(n / \log n)}$ using the Blum--Kalai--Wasserman (BKW) algorithm \cite{BKW03} and its modern refinements \cite{LF06,Lyu05,EKM17,GJL14,ZJW16,WS21}. Grover's algorithm reduces the naive key-search cost from $2^k$ to $2^{k/2}$ \cite{Gro96}, and quantum variants of BKW remain the dominant quantum threat; no published exponential quantum speed-up over the best classical methods is known \cite{Kir11}.

LPN has been the foundation for a wide range of cryptographic primitives, including OWFs, PKE, OT, and MPC \cite{ABW10,Ale11,DGM14,BCGI18}. Its main theoretical limitation is the absence of a worst-case hardness reduction comparable to Regev's LWE reduction from lattice problems \cite{Reg05}. LPN hardness is therefore treated as a \emph{primitive assumption}, accepted on the basis of sustained cryptanalytic resistance rather than reducibility.

Post-quantum cryptography currently rests on six classical hardness families: multivariate quadratic equations (MQ), lattices (LWE/SIS), error-correcting codes (LPN/syndrome decoding), hash functions, Ising models, and tensor problems. A seventh family --- genuinely new, structurally distinct, and not reducible to the existing six --- has been sought since the inception of the field. LSN is the strongest current candidate for such a family.

\subsection{The LSN Problem}
\label{subsec:lsn}

We propose \emph{Lagrangian Subspace Noise} (LSN), also referred to as \emph{Learning Stabilizers with Noise} or \emph{symplectic LPN}. The secret is not a linear function but a Lagrangian subspace $L \subset \F_2^{2n}$ of a symplectic vector space. The learner receives samples
\[
(a_i, b_i) \quad \text{where} \quad b_i = \mathbf{1}_L(a_i) \oplus e_i,
\]
with $a_i \sim \F_2^{2n}$ uniform, $e_i \sim \operatorname{Bernoulli}(p)$, and $\mathbf{1}_L$ the indicator of $L$. The secret space is the Lagrangian Grassmannian:
\begin{equation}
\label{eq:lagr-count}
|\Lagr(2n, \F_2)| = \prod_{i=1}^n (2^i + 1) = 2^{n(n+1)/2 + O(1)}.
\end{equation}
This is exponentially larger than LPN's $2^n$ secrets, yet it is not an unstructured space: every secret satisfies the strong isotropy constraint $L = L^{\perp_\Omega}$. It is exactly this tension --- exponential size versus rigid structure --- that creates both the hardness and the technical difficulty of LSN.

Recent independent work by Khesin \etal\ \cite{KLP+25}, Poremba--Quek--Shor \cite{PQS26}, and Lu \etal\ \cite{LPQR26} has shown that LSN is at least as hard as constant-rate LPN even for a single logical qubit, and has established strong barriers against reducing symplectic LPN to standard LPN in the cryptographically relevant low-noise regime. These results position LSN as the strongest current candidate for a genuinely new post-quantum hardness family outside the six classical ones (multivariate-quadratic, lattice, code, hash, Ising, tensor).

\subsection{Contributions}
\label{subsec:contributions}

Our contributions are organized around proof, barriers, and constructions.

\begin{enumerate}
    \item \textbf{Exact SQ lower bound (\Cref{sec:sq}).} We prove an unconditional $\Omega(2^n)$ SQ query lower bound via an explicit symplectic spread, and a conditional $q \geq 2^{2n - O(1)}$ bound under a pencil-extremality conjecture. Both rest on an exact pairwise-correlation formula
    \[
    \langle D_L, D_{L'} \rangle = \frac{(1-2p)^2}{p(1-p)} \cdot 2^{j-2n},
    \]
    where $j = \dim(L \cap L')$ and $\langle D_L, D_{L'} \rangle$ denotes the pairwise correlation. For $p=1/4$ the coefficient is exactly $4/3$. A Markov concentration argument shows that a subset of size $2^{2n}$ with average correlation $O(2^{-2n})$ exists; the stronger worst-case bound remains open. The results hold for adaptive SQ algorithms.

    \item \textbf{Reduction barriers (\Cref{sec:barriers}).} We prove that linear feature maps give zero advantage against LSN; polynomial feature maps of degree $< n$ have error at least $2^{-n}$; and adaptive linear SQ algorithms are completely blocked by a direct expectation computation. Together these barriers eliminate the natural reduction routes and localize the remaining open question to genuinely non-linear, adaptive reductions.

    \item \textbf{Cryptographic primitives (\Cref{sec:primitives}).}
    \begin{itemize}
        \item \textbf{LSN-SNARK:} A preprocessing SNARK for LSN membership with $O(n^2)$ R1CS constraints via the Siegel-chart (graph-Lagrangian) representation. Verified counts: $64$ for $n=8$, $1{,}764$ for $n=42$, $4{,}225$ for $n=65$.
        \item \textbf{LSN-KEM:} A CCA-secure KEM via concatenated polar codes. Direct polar coding at $p=1/4$ fails ($Z_0 = 0.866$, SC bound $P_e > 1$), so we concatenate an inner repetition code with an outer polar code. For $r=7$ repetition we obtain effective noise $p' = 0.0706$, SC bound $P_e < 2^{-81}$, and parameters PK~=~1.78~KB, CT~=~288~B at 80-bit security. For $r=11$ we obtain $P_e < 2^{-149}$ at 128-bit security.
    \end{itemize}

    \item \textbf{Rigorous caveats (\Cref{sec:open}).} Every claim is classified as theorem, evidence, or conjecture. We state explicitly the standard-model gap --- the absence of a non-vacuous LPN(low-noise) $\to$ LSN or LWE $\to$ LSN reduction --- that must close for full community acceptance.
\end{enumerate}

\section{Preliminaries}
\label{sec:prelim}

\subsection{Symplectic Vector Spaces}
\label{subsec:symplectic}

Let $V = \F_2^{2n}$ with standard symplectic form
\[
\Omega(x,y) = \sum_{i=1}^n (x_{2i-1} y_{2i} + x_{2i} y_{2i-1}).
\]
A subspace $L \subset V$ is \emph{isotropic} if $\Omega|_{L \times L} = 0$, and \emph{Lagrangian} if it is isotropic of dimension $n$; equivalently $L = L^{\perp_\Omega}$. We write $\Lagr(2n, \F_2)$ for the Lagrangian Grassmannian.

\begin{proposition}[Lagrangian count \cite{Mac71}]
$|\Lagr(2n, \F_2)| = \prod_{i=1}^n (2^i + 1) = 2^{n(n+1)/2 + O(1)}$.
\end{proposition}

The intersection of two Lagrangians has even codimension. For $L, L' \in \Lagr(2n)$ let $j = \dim(L \cap L')$. The distribution of $j$ is governed by the $q$-binomial coefficients at $q=2$:
\begin{equation}
\label{eq:j-distribution}
\Pr[j=k] = \frac{{n \brack k}_2 \cdot 2^{(n-k)(n-k+1)/2}}{|\Lagr(2n,\F_2)|}.
\end{equation}

\subsection{Fourier Transform on the Symplectic Space}
\label{subsec:fourier}

The unnormalized $\Omega$-Fourier transform of $f \colon \F_2^{2n} \to \R$ is
\[
\mathcal{F}_\Omega[f](\xi) = \sum_{x \in \F_2^{2n}} f(x) (-1)^{\Omega(x, \xi)}.
\]

\begin{lemma}[Self-duality]
For every $L \in \Lagr(2n,\F_2)$,
\[
\mathcal{F}_\Omega[\mathbf{1}_L] = 2^n \cdot \mathbf{1}_L.
\]
\end{lemma}

\begin{proof}
If $\xi \in L$, then $\Omega(x,\xi)=0$ for all $x \in L$, so the sum is $|L|=2^n$. If $\xi \notin L$, choose $x_0 \in L$ with $\Omega(x_0,\xi)=1$; the involution $x \mapsto x+x_0$ pairs terms of opposite sign, so the sum is $0$.
\end{proof}

\subsection{Statistical Query Model}
\label{subsec:sq}

Let $\mathcal{D}$ be a class of distributions and $D_0$ a reference distribution. The \emph{statistical dimension with average correlation} \cite{FGR+17} is
\[
\SDA(B(\mathcal{D},D_0),\gamma) = \max\Bigl\{d : \forall S \subseteq \mathcal{D}, |S| \geq \frac{|\mathcal{D}|}{d} \Rightarrow \frac{1}{|S|^2}\sum_{D,D' \in S} |\langle D,D' \rangle| \leq \gamma\Bigr\}.
\]

\begin{theorem}[Feldman \etal\ \cite{FGR+17}]
\label{thm:feldman}
If $\SDA(B(\mathcal{D},D_0),\gamma) = d$, any SQ algorithm distinguishing $\mathcal{D}$ from $D_0$ with success probability $\alpha > 1/2$ requires $q \geq (2\alpha - 1)d$ queries to $\VSTAT(1/(3\gamma))$.
\end{theorem}

The theorem applies to randomized \emph{adaptive} SQ algorithms, so our lower bound will inherit adaptivity automatically.

\subsection{Notation}
\label{subsec:notation}

We write $D_L$ for the LSN distribution with secret $L$ and noise rate $p$, and $D_0$ for the noise-only distribution in which $b \sim \operatorname{Bernoulli}(p)$ independently of $a$. All logarithms are base $2$ unless stated otherwise.


\section{Related Work}
\label{sec:related}

\paragraph{LPN foundations and cryptanalysis.}
LPN was introduced by Blum, Furst, Kearns, and Lipton \cite{BFKL94} as a foundation for cryptographic primitives. The dominant classical attack is BKW \cite{BKW03}, later optimized with fast Walsh--Hadamard transforms \cite{LF06}, covering codes \cite{GJL14,ZJW16}, and hybrid information-set-decoding techniques \cite{EKM17,WS21}. Lyubashevsky \cite{Lyu05} gave a polynomial-sample variant, and a comprehensive unified attack framework was developed by Esser, K{\"u}bler, and May \cite{EKM17}. Quantum speed-ups are limited to Grover \cite{Gro96} and quantum BKW \cite{Kir11}; no published exponential improvement is known.

\paragraph{LPN variants.}
Many structured variants have been proposed to improve efficiency: Ring-LPN \cite{HKL+12}, LPN with structured noise \cite{DP12}, and the Learning with Rounding (LWR) problem \cite{BPR12}. All of these retain a secret space of size $2^{O(n)}$. LSN differs fundamentally in that its secret space has size $2^{n^2/2 + O(n)}$ while remaining polynomial-time sampleable via symplectic linear algebra.

\paragraph{Symplectic and quantum coding.}
Stabilizer codes, the quantum analogue of classical linear codes, are naturally described by Lagrangian subspaces of $\F_2^{2n}$ \cite{Got97,CRSS98,NC00}. Quantum state tomography of stabilizer states under depolarizing noise is a closely related learning problem \cite{GL22}. The recent independent work of Khesin \etal\ \cite{KLP+25} proved that constant-rate LPN reduces to LSN even for $k=1$, while Lu \etal\ \cite{LPQR26} developed scrambling techniques for symplectic linear spaces and gave strong evidence that sympLPN does not reduce to LPN in the low-noise regime.

\paragraph{Statistical query hardness.}
The SQ model was introduced by Kearns \cite{Kea98} and developed by Blum \etal\ \cite{BFJ+94}. Feldman, Grigorescu, Reyzin, and Vempala \cite{FGR+17} gave the general SDA framework we use, and Diakonikolas \etal\ \cite{DKS17} applied it to high-dimensional robust estimation. Quantum algorithms that access only classical samples can achieve at most a quadratic Grover speed-up over classical algorithms. Since the SQ lower bound is exponential in $n$, such a speed-up does not affect the asymptotic security parameter, making SQ lower bounds relevant evidence for quantum hardness.

\paragraph{Code-based cryptography and polar codes.}
Code-based cryptography dates to McEliece \cite{McE78} and Prange \cite{Pra62}. Modern standards include Classic McEliece \cite{ABD+21}, HQC \cite{AAC+22}, and BIKE. Polar codes, introduced by Arikan \cite{Ari09}, are capacity-achieving with explicit construction; list decoding was introduced by Tal and Vardy \cite{TV15}, and finite-length analysis was developed by Mori--Tanaka \cite{MT09}, Trifonov \cite{Tri12}, and Hassani \etal\ \cite{HVU14}. We use Bhattacharyya bounds for conservative parameter selection \cite{Sei15}.

\paragraph{Post-quantum KEMs.}
The NIST PQC standardization process selected Kyber (lattice) \cite{SAB+22}, Classic McEliece, and HQC \cite{AAC+22} as KEM candidates. LSN-KEM offers competitive sizes with a fundamentally different assumption.

\section{The LSN Problem}
\label{sec:lsn-problem}

\begin{definition}[LSN$_{n,m,p}$]
\label{def:lsn}
Let $L \subset \F_2^{2n}$ be uniformly random Lagrangian. The LSN distribution $D_L$ over $\F_2^{2n} \times \F_2$ is generated by sampling $a \sim \F_2^{2n}$, $e \sim \operatorname{Bernoulli}(p)$, and outputting $(a, b)$ where $b = \mathbf{1}_L(a) \oplus e$. Given $m$ independent samples from $D_L$, recover $L$.
\end{definition}

\paragraph{Search vs.~decision.} All hardness results in this paper are stated for the \emph{decisional} problem: distinguishing $D_L$ from the noise-only distribution $D_0$ (where $b \sim \operatorname{Bernoulli}(p)$ independently of $a$). The search-to-decision reduction follows the same template as for LPN~\cite{BFKL94}: given a decision oracle and a candidate Lagrangian $L'$, one can verify whether $L'=L$ by testing whether the samples are consistent with $D_{L'}$; a standard hybrid argument then yields search hardness from decision hardness. We therefore treat decisional LSN as the primitive assumption.

\subsection{Distance Distribution and Expected Intersection}
\label{subsec:distance}

The correlation between two LSN distributions depends only on the dimension $j = \dim(L \cap L')$ of the intersection of their secrets. The distribution of $j$ is:

\begin{theorem}[Distance distribution]
\label{thm:distance}
For uniform $L, L' \in \Lagr(2n, \F_2)$,
\[
\Pr[j=k] = \frac{{n \brack k}_2 \cdot 2^{(n-k)(n-k+1)/2}}{|\Lagr(2n,\F_2)|}.
\]
Numerical evaluation shows that $\E[j] \to 0.76$ and $\operatorname{Var}(j) \to 0.60$ as $n \to \infty$.
\end{theorem}

The rapid decay of \Cref{eq:j-distribution} means that two random Lagrangians typically intersect in dimension $0$ or $1$; intersections of dimension $\geq 14$ occur with probability $< 2^{-100}$ for $n \geq 42$.

\subsection{Parameters}
\label{subsec:parameters}

\Cref{tab:parameters} gives the exact security parameters derived in \Cref{sec:sq}.

\begin{table}[ht]
\centering
\caption{LSN security parameters for $p=1/4$.}
\label{tab:parameters}
\begin{tabular}{@{}ccccc@{}}
\hline
Security & $n$ & $\log_2(q_{\min})^{\text{(a)}}$ & $|\Lagr|$ & PK (KEM) \\
\hline
80-bit  & 41  & 80.6  & $\sim 2^{861}$  & 1.78 KB \\
128-bit & 65  & 128.6 & $\sim 2^{2145}$ & 2.78 KB \\
192-bit & 97  & 192.6 & $\sim 2^{4753}$ & --- \\
256-bit & 129 & 256.6 & $\sim 2^{8385}$ & --- \\
\hline
\end{tabular}
\vspace{4pt}
\noindent {\footnotesize $^{\text{(a)}}$ The $q_{\min}$ figures use the conditional bound of \Cref{thm:main-sq-cond}; the unconditional spread bound (\Cref{thm:main-sq-uncond}) gives $\Omega(2^n)$.}
\end{table}

The $q_{\min}$ figures use the exact pairwise-correlation formula of Lemma~\ref{lem:exact-corr}, so the $80$-bit line is met \emph{exactly} at $n=41$ with no asymptotic slack. The constant factor in the exponent is $\approx 0.67$ (below $1$), so parameters with $2n < \lambda$ are not advised without a rigorous upper bound. PK sizes are shown only for the two KEM parameter sets derived in \Cref{sec:kem}; 192/256-bit KEM parameters depend on the inner repetition length $r$ and are left as standard extrapolations.

\section{Decoder Landscape}
\label{sec:decoders}

We review the natural decoder families and explain why each fails against LSN at constant noise.

\paragraph{Linear decoders.}
A linear query has the form $q(a,b) = \langle w, a \rangle + c \cdot b$. Its expectation under $D_L$ is
\[
\E_{D_L}[q] = c \cdot \Pr[b=1] = c \cdot \bigl(p + 2^{-n}(1-2p)\bigr),
\]
which is independent of $L$. Hence linear decoders have exactly zero advantage.

\paragraph{Polynomial decoders.}
The indicator $\mathbf{1}_L$ has an exact degree-$n$ polynomial representation
\[
\mathbf{1}_L(x) = \prod_{j=1}^n (1 + \langle w_j, x \rangle)
\]
where $\{w_j\}$ is a basis for the annihilator of $L$. Any polynomial of degree $< n$ misses at least one factor and therefore has error $\geq 2^{-n}$.

\paragraph{List decoding.}
The list size needed to enumerate all candidate Lagrangians is $|\Lagr| = 2^{n^2/2 + O(n)}$. Even with pruning, list decoding is infeasible for $n \geq 41$.

\paragraph{BKW and ISD.}
BKW would require treating the secret as an $n^2$-bit object, giving complexity $2^{\Omega(n^2 / \log n)}$. Information-set decoding does not apply directly because the code ensemble (Lagrangian indicators) is highly structured rather than random linear.

\paragraph{Structural attacks.}
For very low noise ($p = o(1/n)$), the span of the positive samples reveals $L$. At constant noise $p=1/4$, \cite{LPQR26} and our experiments show that structural attacks are ineffective: the span of positives has full dimension after $O(n)$ samples.

\paragraph{Empirical sample-complexity lower bound.}
We ran brute-force experiments for $n=3,4,5$ using transvection-orbit enumeration of the full Lagrangian family.  In each trial we sampled $L\leftarrow\Lagr$ uniformly, drew $m$ noisy samples from $\LSN_{L,1/4}$, and let a brute-force decoder enumerate all Lagrangians, picking the one with the highest agreement.  The recovery success rates (200 trials for $n=3$, 50 for $n=4$, 20 for $n=5$) were:
\begin{center}
\begin{tabular}{@{}cccccc@{}}
\toprule
$n$ & $|\Lagr|$ & $m=2^{2n-2}$ & $m=2^{2n-1}$ & $m=2^{2n}$ & $m=2\cdot2^{2n}$ \\
\midrule
3 & 135 & $11.5\%$ & $18.5\%$ & $46.5\%$ & $80.0\%$ \\
4 & 2{,}295 & $6.0\%$ & $24.0\%$ & $62.0\%$ & $98.0\%$ \\
5 & 75{,}735 & $5.0\%$ & $20.0\%$ & $75.0\%$ & $100\%$~(20/20) \\
\bottomrule
\end{tabular}
\end{center}
Two trends are visible.  First, the success rate at the floor $m=2^{2n}$ \emph{increases} with $n$ ($46.5\%$, $62\%$, $75\%$), consistent with $2^{2n}$ being the asymptotic threshold (the constant factor in front appears to be below $1$ and to matter less as $n$ grows).  Second, halving the budget to $m=2^{2n-1}$ already drops recovery below $25\%$ even at $n=5$, so the threshold sits at the $2^{2n}$ scale rather than a constant factor below it.  These experiments \emph{empirically support} a $2^{2n}$ sample-complexity scale for the brute-force ML decoder, matching the pairwise floor of Lemma~\ref{lem:exact-corr}; they do not by themselves rule out a more sample-efficient decoder---the rigorous lower bound is the SQ query bound of \Cref{sec:sq}.

\section{Statistical Query Lower Bound}
\label{sec:sq}

This section contains the main technical result: an exponential SQ lower bound with an exact correlation formula.

\subsection{Exact Pairwise Correlation}
\label{subsec:exact-corr}

Let $D_0$ be the noise-only distribution. The likelihood ratio of $D_L$ with respect to $D_0$ is
\[
\ell_L(a,b) = \frac{dD_L}{dD_0}(a,b) - 1
            = \mathbf{1}_L(a) \cdot \beta \cdot \Bigl(\mathbf{1}_{b=1} - \mathbf{1}_{b=0} \frac{p}{1-p}\Bigr),
\]
where $\beta = (1-2p)/p$.

\begin{lemma}[Exact pairwise correlation]
\label{lem:exact-corr}
For $L,L' \in \Lagr(2n,\F_2)$ with $j = \dim(L \cap L')$,
\begin{equation}
\label{eq:exact-corr}
\langle D_L, D_{L'} \rangle
= \frac{(1-2p)^2}{p(1-p)} \cdot 2^{j-2n}.
\end{equation}
\end{lemma}

\begin{proof}
Because $a$ and $b$ are independent under $D_0$, the inner product factors:
\[
\langle D_L, D_{L'} \rangle
= \E_a[\mathbf{1}_L(a) \mathbf{1}_{L'}(a)] \cdot
  \E_b\Bigl[\beta^2 \Bigl(\mathbf{1}_{b=1} - \mathbf{1}_{b=0}\frac{p}{1-p}\Bigr)^2\Bigr].
\]
The first factor is $|L \cap L'| / 2^{2n} = 2^{j-2n}$. The second factor is
\[
\beta^2 \Bigl(p + (1-p)\frac{p^2}{(1-p)^2}\Bigr)
= \frac{(1-2p)^2}{p^2} \cdot \frac{p}{1-p}
= \frac{(1-2p)^2}{p(1-p)}.
\]
Multiplying yields \Cref{eq:exact-corr}. For $p=1/4$ the coefficient is exactly $4/3$.
\end{proof}

\subsection{Average Correlation}
\label{subsec:avg-corr}

Taking expectation over $j$ gives the average pairwise correlation.

\begin{lemma}[Average correlation]
Let $C_n = \E[2^j]$ computed from \Cref{thm:distance}. Then $C_n \to 2$ and
\[
\rho_{\mathrm{avg}}
= \frac{(1-2p)^2}{p(1-p)} \cdot C_n \cdot 2^{-2n}.
\]
\end{lemma}

\subsection{SDA Concentration}
\label{subsec:sda}

\begin{lemma}[Existence of low-correlation subset]
\label{lem:sda}
Let $\gamma = 2\rho_{\mathrm{avg}}$. There exists a subset $\mathcal{D}' \subset \{D_L\}$ of size $|\mathcal{D}'| = 2^{2n}$ with average correlation at most $\gamma$.
\end{lemma}

\begin{proof}
Consider a uniformly random subset $S \subset \Lagr(2n)$ of size $M = 2^{2n}$. By symmetry of the Lagrangian Grassmannian under $\mathrm{Sp}(2n)$, $\E[\rho(S)] = \rho_{\mathrm{avg}}$. Markov's inequality gives
\[
\Pr[\rho(S) > 2\rho_{\mathrm{avg}}] < \frac{\E[\rho(S)]}{2\rho_{\mathrm{avg}}} = \frac{1}{2}.
\]
Hence with probability at least $1/2$, a random subset of size $2^{2n}$ has average correlation at most $\gamma = 2\rho_{\mathrm{avg}}$. This proves existence of such a subset. A worst-case bound over \emph{all} subsets of size $2^{2n}$ would require a structural argument; we discuss this in \Cref{sec:open}.
\end{proof}

\subsection{Unconditional SQ Bound: The Symplectic Spread}
\label{subsec:spread-bound}

We first give an unconditional exponential lower bound on a concrete worst-case promise problem. The bound is weaker in both query count and VSTAT strength than the conditional average-case bound of \Cref{thm:main-sq-cond}, but it requires no unproven structural hypothesis.

Let $V = \F_{2^n} \times \F_{2^n}$ with symplectic form $\omega((a,b),(c,d)) = \operatorname{Tr}(ad) + \operatorname{Tr}(bc)$, where $\operatorname{Tr}$ is the trace from $\F_{2^n}$ to $\F_2$. The \emph{Desarguesian symplectic spread} $\mathcal{S}^* = \{L_\lambda\}_{\lambda \in \F_{2^n}} \cup \{L_\infty\}$ consists of $d_0 = 2^n+1$ Lagrangians
\[
L_\lambda = \{(x, \lambda x) : x \in \F_{2^n}\}, \qquad L_\infty = \{0\} \times \F_{2^n},
\]
which are pairwise transversal: $\dim(L \cap L') = 0$ for all distinct $L, L' \in \mathcal{S}^*$.

\begin{theorem}[Unconditional spread bound]
\label{thm:main-sq-uncond}
Let $\kappa = (1-2p)^2/(p(1-p))$ and fix $1 \leq t \leq n-1$. Set $\bar\gamma_t = \kappa \cdot 2^{-(2n-t)}$. Every subset $T \subseteq \mathcal{S}^*$ with $|T| \geq 2^{n-t}$ has diagonal-inclusive average correlation at most $2\bar\gamma_t$. Consequently $\SDA(\mathcal{S}^*, 2\bar\gamma_t) \geq 2^t$, and any SQ algorithm that is correct on every $L \in \mathcal{S}^*$ requires $\Omega(2^t)$ queries to $\VSTAT(1/(6\bar\gamma_t))$.
\end{theorem}

At $t = n-1$ this yields an unconditional $\Omega(2^n)$ query lower bound at VSTAT strength $O(2^n)$.

\begin{proof}
For distinct $L, L' \in \mathcal{S}^*$ we have $j = \dim(L \cap L') = 0$, so by Lemma~\ref{lem:exact-corr} the pairwise correlation is exactly $\langle D_L, D_{L'}\rangle = \kappa \cdot 2^{-2n} = \bar\gamma_t \cdot 2^{-t} \leq \bar\gamma_t$. The diagonal terms contribute $\beta = \kappa \cdot 2^{-n}$ each, so for $|T| \geq 2^{n-t}$
\[
\frac{1}{|T|^2}\sum_{D,D' \in T} |\langle D,D'\rangle|
\;\leq\; \frac{|T|^2 \cdot \bar\gamma_t + |T| \cdot \beta}{|T|^2}
\;\leq\; \bar\gamma_t + \frac{\beta}{2^{n-t}}
\;\leq\; 2\bar\gamma_t.
\]
Hence $\SDA(\mathcal{S}^*, 2\bar\gamma_t) \geq |\mathcal{S}^*| / 2^{n-t} \geq 2^t$. Applying \Cref{thm:feldman} with $\alpha = 2/3$ gives the query bound.
\end{proof}

\paragraph{Honesty notes.}
(i) \textbf{Worst-case promise only.} The spread has measure $|\mathcal{S}^*|/|\Lagr(2n,\F_2)| \sim 2^n / 2^{n(n+1)/2} \to 0$, so Theorem~\ref{thm:main-sq-uncond} does \emph{not} cover the average-case (uniform $L$) problem; it protects only the promise problem ``$L \in \mathcal{S}^*$?''.
(ii) \textbf{VSTAT trade-off.} The unconditional bound is $\Omega(2^n)$ queries at $\VSTAT(O(2^n))$; the conditional bound below is $2^{2n-O(1)}$ queries at $\VSTAT(\Theta(2^{2n}))$. The two parameters are not interchangeable.
(iii) The subfamily restriction is valid for the promise problem: any algorithm solving LSN over the full family is correct on $\mathcal{S}^*$ a fortiori.

\subsection{Conditional SQ Bound}
\label{subsec:main-sq}

The stronger $2^{2n}$-scale bound rests on a structural conjecture about extremal subsets.

\begin{conjecture}[Pencil extremality]
\label{conj:pencil}
There is an absolute constant $c$ such that every subset $\mathcal{D}' \subseteq \{D_L\}$ of size $|\mathcal{D}'| \geq |\Lagr(2n,\F_2)| / 2^{2n-c}$ has average correlation at most $5\rho_{\mathrm{avg}}$.
\end{conjecture}

Motivation: $k=2$ isotropic pencils force any threshold above $4\rho_{\mathrm{avg}}$ (their expected $2^{\dim(L\cap L')}$ ratio tends to $4$ from below), while $k \geq 3$ pencils fall below the size threshold $|\Lagr|/2^{2n}$ for $n \geq 4$, and mixtures of pencils with random fill dilute quadratically. A proof would require classifying near-extremal subsets.

\begin{theorem}[Conditional main SQ lower bound]
\label{thm:main-sq-cond}
Assume \Cref{conj:pencil}. Any SQ algorithm distinguishing LSN (uniform $L$) from $D_0$ with probability $> 2/3$ requires $q \geq 2^{2n - O(1)}$ queries to $\VSTAT(1/(15\rho_{\mathrm{avg}}))$.
\end{theorem}

\begin{proof}
Under \Cref{conj:pencil} with $c=0$, every subset of size at least $|\Lagr(2n,\F_2)| / 2^{2n}$ has average correlation at most $5\rho_{\mathrm{avg}}$. Hence $\SDA(B(\{D_L\},D_0), 5\rho_{\mathrm{avg}}) \geq 2^{2n}$. Applying \Cref{thm:feldman} with $\alpha = 2/3$ and $\gamma = 5\rho_{\mathrm{avg}}$ yields $q \geq (4/3 - 1) \cdot 2^{2n} = 2^{2n - O(1)}$ queries to $\VSTAT(1/(15\rho_{\mathrm{avg}}))$.
\end{proof}

\begin{theorem}[Adaptive SQ]
\Cref{thm:main-sq-cond} holds for randomized adaptive SQ algorithms.
\end{theorem}

\begin{proof}
\Cref{thm:feldman} applies to randomized adaptive SQ algorithms \cite[Theorem 3.7]{FGR+17}, so the bound is inherited directly.
\end{proof}

\subsection{Adaptive Linear SQ is Blocked}
\label{subsec:adaptive-linear}

\begin{theorem}[Adaptive linear SQ impossibility]
\label{thm:linear-sq}
For any linear query $q(a,b) = \langle w, a \rangle + c \cdot b$, the expectation $\E_{D_L}[q]$ is independent of $L$. Consequently, adaptive linear SQ algorithms have zero advantage.
\end{theorem}

\begin{proof}
Direct computation: $\E_a[\langle w,a \rangle] = 0$ and $\E_{D_L}[b] = p + 2^{-n}(1-2p)$, which does not depend on $L$.
\end{proof}


\section{Quantum Analysis}
\label{sec:quantum}

\subsection{Grover Search}
\label{subsec:grover}

The secret space has size $|\Lagr(2n,\F_2)| = 2^{n^2/2 + O(n)}$. A Grover search over this space requires $2^{n^2/4 + O(n)}$ iterations, each costing at least $\Omega(mn)$ to evaluate an LSN sample set. The total cost $2^{n^2/4 + O(n)}$ dominates all classical bounds for $n \geq 4$, so Grover is not a practical threat.

\subsection{Quantum BKW}
\label{subsec:qbkw}

A quantum version of BKW can speed up the inner Walsh--Hadamard transform and the search for low-weight parity checks, yielding complexity $2^{O(n^2 / \log n)}$ \cite{Kir11}. Since $n^2/\log n \gg 2n$ for all $n \geq 4$, this exceeds the SQ lower bound and is irrelevant for security parameter selection.

\subsection{Quantum SQ}
\label{subsec:quantum-sq}

Because LSN samples are classical bits, any quantum algorithm must either access them through classical queries or perform a quantum walk over the sample space. A Grover search over the $2^{n^2/2}$-sized secret space still requires $2^{n^2/4}$ iterations (\Cref{subsec:grover}), far above the security parameter. Thus the exponential SQ lower bound provides evidence (not proof) that no efficient quantum algorithm solves LSN.

\begin{conjecture}[Quantum LSN]
\label{conj:quantum}
Any quantum algorithm solving LSN$_{n,m,p}$ with constant $p$ requires $\Omega(2^n)$ queries.
\end{conjecture}

\subsection{Dihedral Hidden Subgroup}
\label{subsec:dhsp}

The Lagrangian Grassmannian is not a group; it lacks a transitive abelian structure. Consequently, no reduction to the dihedral hidden-subgroup problem or its generalizations is known.

\subsection{What Is and Is Not Covered}
\label{subsec:quantum-covered}

We summarize the quantum attack landscape honestly.

\paragraph{Blocked vectors.} The following quantum approaches have been analyzed and ruled out for LSN at constant noise:
\begin{itemize}
\item \textbf{Hidden Subgroup Problem.} The Lagrangian Grassmannian lacks a group structure with coset decomposition; no reduction to DHSP or abelian HSP is known.
\item \textbf{Grover search.} Complexity $2^{n^2/2 + O(n)}$ exceeds all classical bounds.
\item \textbf{Quantum BKW.} Complexity $2^{O(n^2 / \log n)}$ is super-exponential in the security parameter.
\end{itemize}

\paragraph{Not covered.} We do \emph{not} claim a quantum lower bound for LSN. Quantum hardness is treated as a \textbf{conjecture} (\Cref{conj:quantum}), analogous to the status of LWE and LPN. A rigorous quantum query lower bound would require new techniques beyond the SQ model and is left as future work.

\section{Cryptographic Primitives}
\label{sec:primitives}

\subsection{LSN-SNARK}
\label{sec:snark}

We construct a preprocessing SNARK for the relation ``$x \in L$'' given a committed Lagrangian $L$. The construction exploits the \emph{Siegel-chart} (graph-Lagrangian) representation to reduce the R1CS circuit from $O(n^3)$ to $O(n^2)$ constraints.

\paragraph{Siegel-chart representation.}
Fix a public symplectic matrix $S \in \mathrm{Sp}(2n,\F_2)$ (e.g., generated from a standard seed). Every ``generic'' Lagrangian can be written as
\[
L_M = S \cdot \{(x, Mx) : x \in \F_2^n\},
\]
where $M \in \mathrm{Sym}(n,\F_2)$ is a symmetric matrix. The set of graph Lagrangians has cardinality $|\mathrm{Sym}(n)| = 2^{n(n+1)/2}$, which differs from the full Lagrangian family $|\Lagr| = \prod_{i=1}^n (2^i+1)$ by only a constant factor ($\approx 2.38$). Consequently, restricting to graph Lagrangians does not affect asymptotic security; the brute-force complexity remains $2^{n^2/2+O(n)}$ and the SQ lower bounds (Theorems~\ref{thm:main-sq-uncond} and~\ref{thm:main-sq-cond}) apply asymptotically unchanged because correlation structure is preserved under symplectic change of coordinates.

\paragraph{Binding $M$ to the public key.}
The LSN-KEM public key (\cref{sec:kem}) contains $m$ noisy samples $(x_i, y_i)$ generated from the secret Lagrangian $L_M$. To prove knowledge of $M$ in a SNARK, the prover must bind their witness to the public key. The public key is therefore augmented with a commitment $c = \mathsf{Hash}(M)$, where $\mathsf{Hash}$ is a SNARK-friendly hash function (e.g., Poseidon). This introduces preimage-resistance of $\mathsf{Hash}$ as a second assumption; we size $\mathsf{Hash}$ ($\geq 256$-bit output, Poseidon) so that quantum preimage search ($2^{128}$) is at least the LSN level, making LSN the bottleneck. The SNARK circuit then verifies two things: (1)~commitment opening ($c = \mathsf{Hash}(M)$), and (2)~membership ($x \in L_M$). Hashing the $n(n+1)/2$ bits of $M$ with Poseidon adds only $O(n)$ constraints---packing the bits into $\approx 9$ field elements for $n=65$ costs roughly $300$--$500$ R1CS constraints. The full circuit size is therefore $n^2 + O(n)$, still $O(n^2)$ and still orders of magnitude below ZK-lattice circuits.

\paragraph{Membership circuit.}
The core membership sub-circuit verifies $x \in L_M$ in two steps:
\begin{enumerate}
\item Apply the public inverse $S^{-1}$ to $x$ to obtain $(a,b) = S^{-1}x$. This is a linear operation over $\F_2$, hence costs \textbf{zero} R1CS constraints.
\item Check $b = Ma$. Each entry $b_i = \sum_{j=1}^n M_{ij} a_j$ is a dot product of length $n$ over $\F_2$. Treating $a$ as a witness vector yields $n$ AND gates per row, for $n^2$ constraints total; the final sums are linear and absorbed into the R1CS output side. Symmetry of $M$ is free: we encode only the upper-triangular $n(n+1)/2$ entries.

\paragraph{Further optimization.} If $a$ is provided directly as a public input (derived from $x$ by the verifier), each product $M_{ij}a_j$ becomes a selector---free when $a_j=0$ and a copy when $a_j=1$. The check $b_i = \sum_j M_{ij}a_j$ then collapses to a single linear constraint per row, giving \textbf{$n$ constraints} over $\F_2$ R1CS, or $n(n+3)/2$ constraints over a prime field with boolean witnesses for $M$.
\end{enumerate}

\Cref{tab:r1cs} compares LSN-SNARK with other ZK-PQC approaches.

\begin{table}[ht]
\centering
\caption{R1CS constraints for LSN membership vs.~ZK-PQC full primitives.}
\label{tab:r1cs}
{\small\begin{tabular}{@{}lcc@{}}
\hline
Scheme & Constraints ($n=65$) & Asymptotic \\
\hline
\textbf{LSN-SNARK membership} & $\mathbf{\approx 4{,}225}$ & $\mathbf{O(n^2)}$ \\
\textbf{LSN-SNARK full ($+\,$hash)} & $\mathbf{\approx 4{,}500}$--$\mathbf{4{,}700}$ & $\mathbf{O(n^2)}$ \\
SHA-256 (per block) & $\approx 25{,}000$ & $O(1)$ \\
AES-128 & $\approx 16{,}000$ & $O(1)$ \\
ZK-Kyber (full KEM verif.) & millions & $O(n \log n)$ \\
ZK-SPHINCS+ (full sig. verif.) & millions & $O(\lambda \log \lambda)$ \\
\hline
\end{tabular}}
\end{table}

\paragraph{Witness structure.} The witness $\mathbf{w}$ consists of:
\begin{itemize}
\item $n(n+1)/2$ bits for the upper-triangular entries of $M \in \mathrm{Sym}(n,\F_2)$;
\item $n$ bits for the vector $a$ (the first half of $S^{-1}x$, computed as a linear function of $x$ but included in the witness for the dot-product checks);
\item auxiliary opening data for the hash commitment (omitted for deterministic hashes, negligible otherwise).
\end{itemize}
Total witness size: $n(n+1)/2 + O(n) = O(n^2)$ bits.

\paragraph{Verified counts.}
The $n^2$ prediction is exact for the membership sub-circuit when $a$ is a witness variable: $n=8$ yields $64$ constraints, $n=41$ (80-bit) yields $1{,}681$, and $n=65$ (128-bit) yields $4{,}225$. With public $a$, the optimized membership counts are $8$, $41$, and $65$ respectively over $\F_2$. The \emph{full} circuit (membership $+$ Poseidon commitment opening) at $n=65$ is $\approx 4{,}500$--$4{,}700$ constraints depending on hash parameterization.

\paragraph{Why this matters.} The membership relation $x \in L_M$ thus has an unusually compact ZK circuit---$\approx 4{,}225$ R1CS constraints at $n=65$ (or $65$ with public $a$), small relative to hashing one SHA-256 block ($\approx 25{,}000$). The full primitive (membership $+$ commitment opening) stays below $\approx 5{,}000$ constraints---about $5\times$ smaller than a single SHA-256 block ($\approx 25{,}000$) and roughly $1{,}000\times$ smaller than published ZK-Kyber circuits (millions). \textbf{Scope.} The table compares \emph{core-relation} circuits; full signature/KEM verification includes additional binding overhead that is comparable across schemes. Our claim is not a record for the full signature, but that the \emph{membership} sub-circuit---the inner relation being proved---is exceptionally compact. That compactness helps building blocks such as \emph{verifiable PQ encryption} and \emph{recursive SNARK aggregation}, where many membership checks dominate the cost.

\paragraph{Batch verification.}
The compactness enables dramatic amortization when proving membership for $k$ distinct elements $x^{(1)},\dots,x^{(k)}$ under the same secret $M$. The commitment $c=\mathsf{Hash}(M)$ is verified \textbf{once}; each additional membership check adds only $n^2$ constraints (or $n$ with public $a$). For batch size $k$:
\[
\text{Total constraints} = \underbrace{O(n)}_{\text{hash opening}} + k \cdot \underbrace{O(n^2)}_{\text{membership}}.
\]
In contrast, ZK-Kyber verification requires independent NTT evaluations per proof, offering no comparable amortization. Concrete example ($n=65$, $k=1000$): LSN uses $\approx 500 + 1000\cdot 4{,}225 \approx 4.2\times 10^6$ constraints, whereas ZK-Kyber at $\approx 5\times 10^6$ constraints per proof would need $\approx 5\times 10^9$---a $1{,}000\times$ gap.

\subsection{LSN-KEM}
\label{sec:kem}

\subsubsection{Why concatenated codes}
\label{subsubsec:why-concat}

Direct polar coding for LSN at $p=1/4$ fails. The Bhattacharyya parameter of BSC($p$) is $Z_0 = 2\sqrt{p(1-p)} = 0.866$, too close to 1 for the finite-length polar code $N=2048$ to achieve any useful reliability. We therefore concatenate:
\begin{itemize}
\item an inner repetition code of length $r$, reducing the effective channel to BSC($p'$) where
  \[
  p' = \sum_{k=\lceil r/2 \rceil + 1}^{r} \binom{r}{k} p^k (1-p)^{r-k};
  \]
\item an outer polar code of length $N=2048$ and dimension $K=256$ over BSC($p'$), decoded with SCL of list size $L=8$.
\end{itemize}

\Cref{tab:kem-params} gives the parameters.

\begin{table}[ht]
\centering
\caption{LSN-KEM parameters for concatenated polar reconciliation.}
\label{tab:kem-params}
\begin{tabular}{@{}cccccc@{}}
\hline
Security & $n$ & $r$ & $p'$ & SC bound & SCL estimate \\
\hline
80-bit  & 41 & 7  & 0.0706 & $2^{-81}$ & $2^{-89}$ \\
128-bit & 65 & 11 & 0.0343 & $2^{-149}$ & $2^{-157}$ \\
\hline
\end{tabular}
\end{table}

\subsubsection{Construction}
\label{subsubsec:kem-construction}

We use the random-oracle model (ROM) for the security proof.

\begin{description}
\item[$\mathsf{KeyGen}(1^\lambda)$:] Let $S \in \mathrm{Sp}(2n,\F_2)$ be a public symplectic matrix. Choose a symmetric matrix $M \leftarrow \mathrm{Sym}(n,\F_2)$ uniformly and define $L = S \cdot \{(x,Mx) : x \in \F_2^n\}$. Choose $\mathsf{seed} \leftarrow \{0,1\}^\lambda$. For $i=1,\dots,m$ (where $m=N \cdot r$), compute $x_i = \mathsf{PRG}(\mathsf{seed}, i)$ and $y_i = \mathbf{1}_L(x_i) \oplus e_i$. Output $\mathsf{pk} = (\mathsf{seed}, y_1,\dots,y_m)$ and $\mathsf{sk} = M$.

\paragraph{Security remark.} Restricting to graph Lagrangians reduces the key-space size by a constant factor $|\mathrm{Sym}(n)| / |\Lagr| \to \prod_{i\ge1}(1+2^{-i})^{-1}\approx1/2.38$, which does not affect the asymptotic security parameter. The SQ lower bounds (Theorems~\ref{thm:main-sq-uncond} and~\ref{thm:main-sq-cond}) and the empirical brute-force barriers (\Cref{sec:decoders}) apply asymptotically unchanged because correlation structure is preserved under symplectic change of coordinates. As a bonus, the secret key is more compact: $M$ requires only $n(n+1)/2$ bits versus $2n^2$ bits for a general basis representation.

\item[Encaps:] Choose randomness $\rho=(s,u)$ with $s \leftarrow \{0,1\}^\lambda$ and $u \leftarrow \{0,1\}^K$. Compute a Fisher--Yates permutation $\pi = \mathsf{Permutation}(\mathsf{PRG}(s))$ over $[m]$, and let $\mathsf{indices} = \pi[:Nr]$. Group the selected indices into $N$ blocks of length $r$ and compute $v_j = \mathsf{majority}(y \text{ in block } j)$ for each block. Compute $c = \mathsf{PolarEncode}(u)$, set $\mathsf{ct} = (s, v \oplus c)$, and $K = \mathsf{Hash}(s, u)$.

\item[Decaps:] Recompute $\pi$ and the blocks. For each block compute $v'_j = \mathsf{majority}(\mathbf{1}_L(x))$ using the secret key. Because the $x$'s are derived from the same public seed, the $v'_j$ are deterministic; they need not equal the noisy $v_j$ sent by the encapsulator, but the difference $v_j \oplus v'_j$ is exactly the effective noise seen by the polar decoder. Compute $c' = \mathsf{syn} \oplus v'$, decode $u' = \mathsf{PolarSCLDecode}(c')$, and output $K = \mathsf{Hash}(s, u')$ (or $\perp$ if decoding fails).
\end{description}

The permutation-based index sampling guarantees distinct indices, eliminating birthday collisions. Public-key size is $\lambda + N r$ bits; ciphertext size is $\lambda + N$ bits.

\paragraph{Implementation note.} A small-scale Python prototype of the SCL decoder exhibited a bug (BLER = 1.0 for $N=128,256$), contradicting the conservative Bhattacharyya bound and indicating an implementation issue rather than a parameter problem. Subsequent validation with the verified \texttt{komm} SC decoder \cite{komm} revealed a second bug in the Bhattacharyya parameter indexing: the original code concatenated all bad channels before all good channels, whereas \texttt{komm}'s natural-order polar code requires interleaving ($z_{2i}=2z_i-z_i^2$, $z_{2i+1}=z_i^2$). After correcting the frozen-set construction, all tested configurations yielded BLER $=0$: at $p'=0.0706$ for $N=128,256,512$ with $K=N/8$, and at $p'=0.0343$ for all tested lengths. The previous $N=256$ anomaly (BLER $=0.115$) was entirely due to the indexing mismatch, not incomplete polarization. These short-length, 200-trial simulations are \emph{preliminary}: they certify only $\mathrm{BLER}\lesssim 1/200$, not the design point $2^{-80}$, and they do not test the design length $N=2048$. The conservative Ar\i{}kan--Bhattacharyya bound ($2^{-80}$ at $r=7$, $2^{-148}$ at $r=11$) remains the design basis; full $N=2048$ validation is deferred to the production Rust implementation.

\subsubsection{Correctness}
\label{subsubsec:kem-correctness}

\begin{theorem}[Correctness]
\label{thm:kem-correctness}
For $r=7$, $\Pr[K \neq K'] < 2^{-80}$; for $r=11$, $\Pr[K \neq K'] < 2^{-128}$.
\end{theorem}

\begin{proof}
Inner majority vote converts BSC($1/4$) into BSC($p'$). The outer polar code Bhattacharyya bound gives $P_e \leq \frac{1}{2}\sum_{i \in \mathcal{I}} Z(W_i)$, where $\mathcal{I}$ is the information set. Numerical evaluation of the bound yields the claimed values; SCL decoding with $L=8$ improves them by a small constant factor.
\end{proof}

\subsubsection{IND-CPA Security}
\label{subsubsec:kem-indcpa}

\begin{theorem}[IND-CPA]
LSN-KEM is IND-CPA secure in the random-oracle model under the LSN hardness assumption.
\end{theorem}

\begin{proof}[Proof sketch]
The proof uses three game hops.
\begin{enumerate}
\item \textbf{Game 0 $\to$ Game 1:} Replace $\mathsf{PRG}$ with a truly random function. Indistinguishability follows from PRG security.
\item \textbf{Game 1 $\to$ Game 2:} Replace the public-key labels $(y_1,\dots,y_m)$ with uniform random bits. We construct a direct reduction $\mathcal{B}$ from KEM IND-CPA to LSN: $\mathcal{B}$ receives an LSN challenge $(x_i, y_i^*)$, programs $\mathsf{seed}$ so that $\mathsf{PRG}(\mathsf{seed},i)=x_i$, and sets $\mathsf{pk}=(\mathsf{seed}, y_1^*,\dots,y_m^*)$. The KEM challenger computes $\mathsf{ct}^*$ using all $m=Nr$ samples. If the LSN challenge is real, the adversary sees Game 1; if random, Game 2. Any advantage $\varepsilon$ in distinguishing the games translates to advantage $\varepsilon$ against LSN. There is no $1/r$ security loss because the reduction uses the LSN challenge directly.
\item \textbf{Game 2 $\to$ Game 3:} In Game 2 the vector $v$ is uniform and independent of the polar codeword $c=\mathsf{PolarEncode}(u)$. Hence $\mathsf{syn}=v \oplus c$ is uniform (one-time pad), and $K=\mathsf{Hash}(s,u)$ is uniform and independent of $\mathsf{ct}^*$ by the random-oracle property.
\qedhere
\end{enumerate}
\end{proof}

\subsubsection{Sizes and Comparison}
\label{subsubsec:kem-sizes}

\Cref{tab:kem-sizes} compares LSN-KEM with selected post-quantum KEMs.

\begin{table}[ht]
\centering
\caption{KEM size comparison.}
\label{tab:kem-sizes}
\begin{tabular}{@{}lcccc@{}}
\hline
Scheme & $|\mathsf{pk}|$ & $|\mathsf{ct}|$ & Assumption & Failure \\
\hline
Kyber-512  & 800 B  & 768 B  & Module-LWE   & $2^{-164}$ \\
HQC-128    & 2.25 KB & 4.50 KB & Syndrome Dec. & $2^{-128}$ \\
LSN-80     & 1.78 KB & 288 B  & LSN          & $< 2^{-89}$ \\
LSN-128    & 2.78 KB & 288 B  & LSN          & $< 2^{-157}$ \\
\hline
\end{tabular}
\end{table}

\subsubsection{CCA Security}
\label{subsubsec:kem-cca}

We obtain IND-CCA security via the Fujisaki--Okamoto (FO) transform \cite{FO99,HHK17} with implicit rejection. Let $\mathsf{BaseKEM}=(\mathsf{KeyGen},\mathsf{Encaps},\mathsf{Decaps})$ be the LSN-KEM of \cref{subsubsec:kem-construction}; recall that $\mathsf{Encaps}$ is deterministic once the random pair $(s,u)$ is fixed. Define a hash function $H:\{0,1\}^*\to\{0,1\}^\lambda$ (modelled as a random oracle).

\begin{description}
\item[$\mathsf{CCA\text{-}KeyGen}(1^\lambda)$:] $(\mathsf{pk},\mathsf{sk})\gets\mathsf{KeyGen}(1^\lambda)$. Output $\mathsf{sk}'=(\mathsf{sk},d)$ where $d\gets\{0,1\}^\lambda$ is an implicit-rejection secret.

\item[$\mathsf{CCA\text{-}Encaps}(\mathsf{pk})$:] Choose $s\gets\{0,1\}^\lambda$ and $u\gets\{0,1\}^K$, compute $(c_0,k_0)\gets\mathsf{Encaps}(\mathsf{pk};(s,u))$, set $c=c_0$ and $K=H(s,u,c_0)$. Output $(c,K)$.

\item[$\mathsf{CCA\text{-}Decaps}(\mathsf{sk}',c)$:] Parse $c=(s,\mathsf{syn})$. Recompute the blocks from $s$ and $v'_j=\mathsf{majority}(\mathbf{1}_L(x))$ as in base decaps. Compute $c'=\mathsf{syn}\oplus v'$ and decode $u'=\mathsf{PolarSCLDecode}(c')$. If $u'=\perp$, output $K=H(d,c)$ \emph{(implicit rejection)}. Re-encapsulate $(c_0',k_0')\gets\mathsf{Encaps}(\mathsf{pk};(s,u'))$. If $c_0'=c$, output $K=H(s,u',c)$; otherwise output $K=H(d,c)$ \emph{(implicit rejection)}.
\end{description}

\begin{theorem}[IND-CCA]
If $\mathsf{BaseKEM}$ is $\delta$-correct with $\delta\leq 2^{-80}$ and IND-CPA secure, then $\mathsf{CCA\text{-}KEM}$ is IND-CCA secure in the random-oracle model.
\end{theorem}

\begin{proof}[Proof sketch]
The FO transform preserves IND-CPA security and adds CCA robustness via the re-encapsulation check. Implicit rejection ensures that decapsulation-failure events are indistinguishable from pseudorandom outputs; we inherit the multi-challenge ROM bound of \cite{HHK17}. CCA security implies non-malleability, so tampered ciphertexts are rejected by the re-encapsulation check. The concatenated code affects only the correctness error $\delta$; since $\delta\leq 2^{-80}$ (Theorem~\ref{thm:kem-correctness}), the standard FO analysis \cite{HHK17} applies unchanged.
\end{proof}

\subsubsection{Multi-User Security}
\label{subsubsec:kem-multiuser}

In the multi-user setting, $N$ parties generate independent key pairs $(\mathsf{pk}_i,\mathsf{sk}_i)$ with independent secrets $M_i\gets\mathrm{Sym}(n,\F_2)$. The adversary receives all $N$ public keys and may query decapsulation oracles for any user.

\paragraph{Reduction to single-user security.}
All users share only the public symplectic matrix $S$; every other parameter ($\mathsf{seed}_i$, noise $e_{i,j}$, permutation $\pi_i$) is independent. Because the $M_i$ are independent and the PRG seeds are independent, the $N$ instances are \emph{statistically independent} conditioned on $S$. By a standard hybrid argument over the $N$ users, any adversary $\mathcal{A}$ against $N$-user IND-CCA security with advantage $\varepsilon_{N}$ yields an adversary against single-user IND-CCA security with advantage $\varepsilon_{1}\geq\varepsilon_{N}/N$.

\paragraph{Key-reuse analysis.}
A more subtle concern is repeated encapsulation to the \emph{same} user. Over $q$ encapsulations the adversary sees $q\cdot m=qNr$ samples $(x_j,y_j)$ from the same secret $M$. The brute-force decoder must still search $2^{n(n+1)/2}$ symmetric matrices, and the conditional SQ lower bound (Theorem~\ref{thm:main-sq-cond}) requires $q_{\mathsf{sq}}\geq 2^{2n-O(1)}$ queries. Translating sample count to query count via $q_{\mathsf{sq}}=qNr$, the number of encapsulations needed for non-negligible advantage satisfies
\[
q\;\gtrsim\;\frac{2^{2n-O(1)}}{Nr}.
\]
(We treat each public sample as at most one SQ query; the R2a experiment supports a $2^{2n}$ \emph{sample} threshold.) For $n=65$, $N=2048$, $r=11$ we have $Nr\approx 2^{14.5}$, giving $q\gtrsim 2^{114}$---far beyond any realistic threat. We therefore treat multi-user and key-reuse security as \emph{straightforward} consequences of the single-user bound, with no special vulnerability introduced by the LSN structure.

\subsection{Implementation Security}
\label{subsec:impl-sec}

\paragraph{Constant-time by construction.}
All core LSN operations are native $\F_2$ operations---XOR, AND, and bit-permute---which are straightforward to implement in constant time on modern CPUs. There is no modular reduction with an odd modulus, no discrete Gaussian sampling, no rejection sampling, and no variable-time branching on secret data. This makes a constant-time implementation \emph{straightforward} rather than requiring careful engineering as in lattice schemes.

\begin{table}[ht]
\centering
\caption{Implementation-security comparison: Kyber vs.~LSN.}
\begin{tabular}{@{}lcc@{}}
\hline
Operation & Kyber & LSN \\
\hline
Modular reduction & Variable-time ($q=3329$) & \textbf{None} ($\F_2$) \\
NTT butterfly & Twiddle-dependent memory access & \textbf{Fixed} XOR pattern \\
Noise sampling & Discrete Gaussian (rejection) & \textbf{Uniform} bits \\
Decoder & Noisy polynomial rounding & \textbf{Majority} $+$ fixed-tree SC \\
Secret-key format & $2n^2$ bits (general basis) & $\mathbf{n(n+1)/2}$ bits (symmetric $M$) \\
\hline
\end{tabular}
\end{table}

\paragraph{Side-channel surface.}
The only non-linear operation on secret data is the majority vote in the inner repetition decoder. A popcount implementation is data-independent on all modern architectures (ARM NEON, AVX-512 VPOPCNTDQ). The outer polar SC decoder traverses a fixed butterfly tree of depth $\log_2 N$; the memory access pattern depends only on $N$, not on the secret $M$ or the noise. Consequently, timing attacks on the decoder require exploiting micro-architectural effects in the popcount instruction itself---far narrower than the modular-reduction and Gaussian-sampling side channels of lattice KEMs.

\paragraph{Current status.} A Python prototype validates parameters and algorithms. A production constant-time Rust implementation (including constant-time popcount, bit-sliced symplectic transforms, and a constant-time SC decoder) is planned for full $N=2048$ validation and KAT vector generation.

\section{Reduction Barriers and the Standard-Model Gap}
\label{sec:barriers}

\Cref{tab:barriers} summarizes the status of natural reductions.

\begin{table}[ht]
\centering
\caption{Reduction landscape for LSN.}
\label{tab:barriers}
\begin{tabular}{@{}lll@{}}
\hline
Class & Status & Reason \\
\hline
Linear             & BLOCKED  & $\E[q]$ is $L$-independent \\
Polynomial degree $< n$ & BLOCKED  & Error $\geq 2^{-n}$ \\
Adaptive linear SQ & BLOCKED  & Direct computation \\
LPN($k=\Theta(n^2)$)   & VACUOUS  & BKW cost $2^{\Omega(n^2/\log n)}$ \\
Beyond polynomial feature maps & OPEN & No proof or construction known \\
\hline
\end{tabular}
\end{table}

The work of Lu \etal\ \cite{LPQR26} proves a strong barrier against \emph{linear} reductions sympLPN $\to$ LPN via the entropy-deficiency argument of Lu \etal\: isotropy constraints cost $C(n,2)$ bits of entropy, and any linear reduction destroys the error before uniformity is reached. Our \Cref{thm:linear-sq} gives a complementary, direct barrier in the SQ model.

The remaining open question is whether a \emph{genuinely non-linear} or \emph{adaptive} reduction exists. We do not claim impossibility. Instead we observe a win-win structure: if such a reduction is found, it would likely improve random self-reductions for LPN itself, making it valuable even if it collapses LSN to the 6.5th family.

\paragraph{The symplectic structure is reduction-inert.}
A sharper observation comes from the structure of the LSN samples themselves. The deterministic quadratic constraint $S_A = 0$ (the $\Omega$-Gram of the public matrix columns) is \emph{public} and \emph{$x$-free}: it depends only on the public matrix $A$, not on the secret $x$ or the noise. Consequently, $S_A=0$ adds no $x$-information to an attacker and provides no usable handle for a reductionist trying to embed LPN into LSN. This mirrors the accepted status of Ring-LWE, which is regarded as a lattice-family assumption by \emph{source} (ring structure native to the problem) rather than by reduction to plain LWE. We formalize this as:

\begin{conjecture}[Source novelty]
\label{conj:source}
The symplectic structure of LSN --- including self-duality $\mathcal{F}_\Omega[\mathbf{1}_L] = 2^n \mathbf{1}_L$, Lagrangian-subspace secrets, and stabilizer degeneracy --- has no analogue in LPN. This structural distinction is the basis for treating LSN as an independent hardness assumption.
\end{conjecture}

\Cref{conj:source} is not decidable by reduction analysis alone; it rests on the absence of a known reduction and the structural barriers documented above.

\section{Open Problems}
\label{sec:open}

We state precisely the problems whose resolution would close the standard-model gap.

\begin{enumerate}
\item \textbf{Statistical dimension under $S_A=0$ \cite{LPQR26}.} The full SQ proof is blocked by the deterministic quadratic symplectic relation $S_A=0$ (the $\Omega$-Gram of columns). Prove that conditioning on $S_A=0$ does not reduce the statistical dimension below $2^{\Omega(n)}$.

\item \textbf{LPN(low-noise) $\to$ LSN.} In the cryptographically relevant regime $p=O(1/\sqrt{n})$, the reduction of \cite{KLP+25} is vacuous. Does a non-vacuous reduction exist, or can its non-existence be proved?

\item \textbf{LWE $\to$ LSN.} A reduction from LWE (or worst-case lattice problems) to LSN would provide the gold-standard foundation currently missing.

\item \textbf{Quantum lower bound.} Prove \Cref{conj:quantum}: any quantum algorithm solving LSN requires $\Omega(2^n)$ queries.

\item \textbf{Optimal polar-code rate.} Maximize the rate of the outer polar code while keeping decryption failure below $2^{-128}$.

\item \textbf{Pencil-extremality conjecture.} \label{open:sda} Prove \Cref{conj:pencil}: that isotropic pencils ($k \leq 2$) are extremal at scale $|\Lagr(2n,\F_2)|/2^{2n}$, giving average correlation $\leq 5\rho_{\mathrm{avg}}$ for every subset of that size. This would upgrade \Cref{thm:main-sq-cond} from a conditional to an unconditional worst-case bound.
\end{enumerate}

\subsection{Honest Limitations of the Present Work}
\label{subsec:limitations}

We flag five practical gaps that a reader should weigh against the theoretical claims above.

\begin{enumerate}
\item \textbf{Empirical gap at $N=2048$.} The concatenated code is designed via the Bhattacharyya bound, which guarantees $P_e\leq 2^{-80}$ for the target block length $N=2048$. We have validated the SC decoder empirically for $N\in\{128,256,512\}$ (all BLER $=0$ in 200 trials), but a direct Monte-Carlo run at $N=2048$ is computationally expensive and has not been completed. The correctness claim for $N=2048$ therefore rests on the Bhattacharyya bound and the smaller-scale validation, not on a direct measurement.

\item \textbf{No constant-time reference implementation.} While \cref{subsec:impl-sec} argues that LSN-KEM is constant-time ``by construction,'' this property has not been verified on a concrete, audited implementation. A production-grade Rust reference implementation with known-answer test (KAT) vectors is planned but not yet available. Physical fault-injection resistance (e.g., bit-flip attacks on the majority vote) is likewise deferred to the audited implementation.

\item \textbf{Full-protocol SNARK is future work.} \Cref{sec:snark} optimizes only the \emph{membership} sub-circuit ($n^2$ constraints). A complete zero-knowledge proof of correct encapsulation---including the polar encoder, the majority-vote decoder, and the random permutation---remains open. The comparison in Table~\ref{tab:r1cs} should be read as a lower bound on what a full SNARK would cost.

\item \textbf{Loose multi-user reduction.} \Cref{subsubsec:kem-multiuser} uses a hybrid argument losing a factor $N$. A tight reduction (random self-reducibility) is not known for LSN. In practice the margin is comfortable ($q\gtrsim 2^{114}$), but the bound is not tight.

\item \textbf{Practical optimizations deferred.} The canonical noise rate $p=1/4$ is fixed; whether LSN remains viable at other rates is unanalyzed. The inner repetition code uses a majority-vote decoder for constant-time implementability, but MAP decoding would achieve a lower effective noise. No public-key compression analogous to structured LPN is known. These are engineering tuning degrees of freedom, not structural vulnerabilities.
\end{enumerate}

\paragraph{Threshold cryptography.} The LSN secret $M$ admits a natural linear secret-sharing scheme over $\F_2$ ($M = \sum M_j$), and the public samples $(x_i,y_i)$ provide a publicly verifiable consistency check because the linear relation $b = Ma = \sum_j M_j a$ is additively shareable. This enables distributed key generation where $n$ parties jointly sample $M$ without any single party learning it. Formalizing robustness against malicious minorities is future work.

\paragraph{Quantum ECC interoperability.} Every LSN secret $L$ is a stabilizer state ($k=0$). Its minimum stabilizer weight is $\Omega(n)$ with high probability by the quantum Gilbert--Varshamov bound, so the secret carries natural quantum error-correcting structure. Whether this classical--quantum linkage yields concrete protocols (e.g., quantum-authenticated key exchange) is an open application direction.

\appendix

\section{Full Proof of Pairwise Correlation}
\label{app:corr}

Here we give a self-contained derivation of \Cref{lem:exact-corr}. Let $D_0$ be the distribution on $\F_2^{2n} \times \F_2$ in which $a$ is uniform and $b \sim \operatorname{Bernoulli}(p)$ independently. For a fixed Lagrangian $L$, the LSN distribution $D_L$ has density
\[
D_L(a,b) = \frac{1}{2^{2n}} \Bigl( (1-p)\mathbf{1}_{b=\mathbf{1}_L(a)} + p\mathbf{1}_{b \neq \mathbf{1}_L(a)} \Bigr).
\]
The likelihood ratio is
\[
\frac{D_L(a,b)}{D_0(a,b)} = 1 + \mathbf{1}_L(a) \cdot \beta \cdot \Bigl(\mathbf{1}_{b=1} - \mathbf{1}_{b=0}\frac{p}{1-p}\Bigr),
\]
where $\beta = (1-2p)/p$. Therefore
\[
\ell_L(a,b) = \mathbf{1}_L(a) \cdot \beta \cdot \Bigl(\mathbf{1}_{b=1} - \mathbf{1}_{b=0}\frac{p}{1-p}\Bigr).
\]
For two secrets $L,L'$ with $j = \dim(L \cap L')$,
\begin{align*}
\langle D_L, D_{L'} \rangle
&= \E_{D_0}[\ell_L \ell_{L'}] \\
&= \E_a[\mathbf{1}_L(a)\mathbf{1}_{L'}(a)] \cdot
   \beta^2 \E_b\Bigl[\Bigl(\mathbf{1}_{b=1} - \mathbf{1}_{b=0}\frac{p}{1-p}\Bigr)^2\Bigr] \\
&= 2^{j-2n} \cdot \beta^2 \Bigl( p + (1-p)\frac{p^2}{(1-p)^2} \Bigr) \\
&= 2^{j-2n} \cdot \frac{(1-2p)^2}{p(1-p)}.
\end{align*}
Setting $p=1/4$ gives the coefficient $(1/2)^2 / (1/4 \cdot 3/4) = (1/4)/(3/16) = 4/3$ exactly.

\section{Q-Binomial Distance Distribution}
\label{app:qbin}

The number of $k$-dimensional subspaces of $\F_2^n$ is the Gaussian binomial coefficient
\[
{n \brack k}_2 = \prod_{i=0}^{k-1} \frac{2^{n-i}-1}{2^{k-i}-1}.
\]
To count pairs of Lagrangians with intersection dimension $k$, fix a $k$-dimensional isotropic subspace $W$. The number of Lagrangians containing $W$ equals the number of Lagrangians in the symplectic quotient $W^\perp / W \cong \F_2^{2(n-k)}$, namely $\prod_{i=1}^{n-k}(2^i+1)$. Summing over all $k$-dimensional subspaces $W \subseteq L$ and using the identity $\prod_{i=1}^{n-k}(2^i+1) = 2^{(n-k)(n-k+1)/2} \prod_{i=1}^{n-k}(1+2^{-i})$, the exact probability is
\[
\Pr[j=k] = \frac{{n \brack k}_2 \cdot 2^{(n-k)(n-k+1)/2}}{\prod_{i=1}^n (2^i+1)}.
\]
The moment generating function $C_n = \E[2^j]$ converges rapidly to $2$ as $n \to \infty$.

\section{Reduction Barrier Summary}
\label{app:barrier}

\paragraph{Linear.} For any linear classifier $h(a) = \langle w,a \rangle + t$, the advantage
\[
\Adv(h) = \bigl| \Pr_{D_L}[h(a)=1] - \Pr_{D_0}[h(a)=1] \bigr|
\]
is zero because $\Pr_{D_L}[h(a)=1] = \Pr_{D_0}[h(a)=1]$.

\paragraph{Polynomial degree $<n$.} The exact polynomial representation of $\mathbf{1}_L$ has degree exactly $n$ and $2^n$ monomials. Any polynomial of degree $d < n$ cannot vanish on $L^c$, so its correlation with $\mathbf{1}_L$ is at most $2^{-n}$.

\paragraph{Adaptive linear SQ.} For $q(a,b) = \langle w,a \rangle + cb$ we have $\E_{D_L}[q] = c(p + 2^{-n}(1-2p))$, independent of $L$. Since every query's expectation is $L$-independent, no sequence of adaptive linear queries can distinguish $\{D_L\}$ from $D_0$.

\section{Extension to $\F_q$}
\label{app:fq}

All results in the main text are stated over $\F_2$. Here we sketch the extension to arbitrary finite fields $\F_q$, which widens the parameter space and allows a trade-off between block length $n$ and field size $q$.

\subsection{Symplectic Spaces over $\F_q$}
\label{subsec:fq-symplectic}

Let $q$ be a prime power and $V = \F_q^{2n}$ equipped with a non-degenerate alternating bilinear form $\Omega(x,y) = x^T J y$ where $J = \begin{psmallmatrix} 0 & I_n \\ -I_n & 0 \end{psmallmatrix}$. A subspace $L \subset V$ is \emph{Lagrangian} if it is maximal isotropic; then $\dim L = n$ and $|\Lagr(2n,\F_q)| = \prod_{i=1}^n (q^i+1)$.

For fixed $L$ and $j = \dim(L \cap L')$, the number of Lagrangians $L'$ with intersection dimension $j$ is
\[
N_j(n,q) = {n \brack j}_q \cdot q^{(n-j)(n-j+1)/2},
\]
where ${n \brack j}_q$ is the Gaussian binomial coefficient. Summing over $j$ recovers $|\Lagr(2n,\F_q)|$.

\subsection{Exact Correlation over $\F_q$}
\label{subsec:fq-corr}

The label remains binary ($b \in \F_2$) even though the ambient space is $\F_q^{2n}$, so the SQ machinery of Feldman \etal\ applies without modification. The likelihood ratio calculation is identical with $2$ replaced by $q$:

\begin{lemma}[$\F_q$ pairwise correlation]
\label{lem:fq-corr}
For $L,L' \in \Lagr(2n,\F_q)$ with $j = \dim(L \cap L')$,
\[
\langle D_L, D_{L'} \rangle = \frac{(1-2p)^2}{p(1-p)} \cdot \frac{q^j}{q^{2n}}.
\]
\end{lemma}

\begin{lemma}[$\F_q$ average correlation]
\label{lem:fq-avg}
Let $C_{n,q} = \E[q^j]$ computed via the $q$-binomial distribution. Then
\[
\rho_{\mathrm{avg}} = \frac{(1-2p)^2}{p(1-p)} \cdot C_{n,q} \cdot q^{-2n}.
\]
Numerically, $C_{n,q} \to 2$ as $n \to \infty$ for every $q$ (the limit is independent of the field size).
\end{lemma}

\subsection{Security Parameters over $\F_q$}
\label{subsec:fq-params}

The SQ lower bound becomes $q_{\mathrm{queries}} \geq q^{2n-O(1)}$. One ``bit'' of security corresponds to $\log_2(q)$ field operations, so the effective bit-security is $2n \log_2(q)$. \Cref{tab:fq-params} gives the block lengths $n$ needed for target security levels.

\begin{table}[ht]
\centering
\caption{LSN parameters over $\F_q$ for $p=1/4$.}
\label{tab:fq-params}
\begin{tabular}{@{}ccccc@{}}
\hline
Security & $\F_2$: $n$ & $\F_3$: $n$ & $\F_5$: $n$ & $\F_7$: $n$ \\
\hline
80-bit  & 41  & 26 & 18 & 15 \\
128-bit & 65  & 41 & 28 & 24 \\
192-bit & 97  & 62 & 42 & 35 \\
256-bit & 129 & 82 & 56 & 46 \\
\hline
\end{tabular}
\end{table}

Larger $q$ reduces the required $n$ but increases the cost of each sample (field arithmetic in $\F_q$ versus $\F_2$). The optimal choice depends on the target platform.

\subsection{Open Questions}
\label{subsec:fq-open}

\begin{enumerate}
\item Extend to $q$-ary labels $y \in \F_q$ with noise over $\F_q$; this requires $q$-ary SQ lower bounds.
\item The proof holds for $q = p^k$ with $k > 1$ because symplectic theory is field-agnostic.
\item In characteristic $2$, alternating is equivalent to symmetric; in odd characteristic they differ. The intersection distribution formula remains valid in all characteristics.
\end{enumerate}

% Bibliography
\bibliographystyle{alpha}
\begin{thebibliography}{99}

\bibitem[AAC+22]{AAC+22}
C.~Adjiman, N.~Aragon, S.~Bettaieb, L.~Bidoux, O.~Blazy, J.-C.~Deneuville, P.~Gaborit, S.~Ghosh, A.~Hauteville, and G.~Zemor.
HQC.
NIST PQC Round 4 Submission, 2022.

\bibitem[ABD+21]{ABD+21}
M.~Albrecht, D.~Bernstein, T.~Chou, C.~Cid, J.~Gilcher, T.~Lange, V.~Maram, I.~von Maurich, R.~Misoczki, R.~Niederhagen, K.~Paterson, E.~Persichetti, C.~Peters, P.~Schwabe, N.~Sendrier, J.~Szefer, C.~Tjhai, M.~Tomlinson, and W.~Wang.
Classic McEliece.
NIST PQC Round 4 Submission, 2021.

\bibitem[ABW10]{ABW10}
A.~Akavia, S.~Goldwasser, and V.~Vaikuntanathan.
Simultaneous hardcore bits and cryptography against memory attacks.
In \emph{TCC}, pages 474--495. Springer, 2010.

\bibitem[Ale11]{Ale11}
A.~Aleknovich.
More on average case vs approximation complexity.
\emph{Computational Complexity}, 20(4):755--786, 2011.

\bibitem[Ari09]{Ari09}
E.~Arikan.
Channel polarization: A method for constructing capacity-achieving codes for symmetric binary-input memoryless channels.
\emph{IEEE Trans.\ Inform.\ Theory}, 55(7):3051--3073, 2009.

\bibitem[BBB+18]{BBB+18}
B.~B\"unz, J.~Bootle, D.~Boneh, A.~Poelstra, P.~Wuille, and G.~Maxwell.
Bulletproofs: Short proofs for confidential transactions and more.
In \emph{IEEE S\&P}, pages 315--334, 2018.

\bibitem[BCTV14]{BCTV14}
E.~Ben-Sasson, A.~Chiesa, D.~Tromer, and M.~Virza.
Scalable zero knowledge via cycles of elliptic curves.
In \emph{CRYPTO}, pages 276--294. Springer, 2014.

\bibitem[BFJ+94]{BFJ+94}
A.~Blum, M.~Furst, J.~Jackson, M.~Kearns, Y.~Mansour, and S.~Rudich.
Weakly learning DNF and characterizing statistical query learning using Fourier analysis.
In \emph{STOC}, pages 253--262. ACM, 1994.

\bibitem[BFKL94]{BFKL94}
A.~Blum, A.~Furst, M.~Kearns, and R.~Lipton.
Cryptographic primitives based on hard learning problems.
In \emph{CRYPTO}, pages 278--291. Springer, 1994.

\bibitem[BKW03]{BKW03}
A.~Blum, A.~Kalai, and H.~Wasserman.
Noise-tolerant learning, the parity problem, and the statistical query model.
\emph{J.\ ACM}, 50(4):506--519, 2003.

\bibitem[BPR12]{BPR12}
A.~Banerjee, C.~Peikert, and A.~Rosen.
Pseudorandom functions and lattices.
In \emph{EUROCRYPT}, pages 719--737. Springer, 2012.

\bibitem[BCGI18]{BCGI18}
E.~Boyle, G.~Conteau, N.~Gilboa, and Y.~Ishai.
Compressing vector OLE.
In \emph{ACM CCS}, pages 896--912, 2018.

\bibitem[CRSS98]{CRSS98}
A.~R.~Calderbank, E.~M.~Rains, P.~W.~Shor, and N.~J.~A.~Sloane.
Quantum error correction via codes over GF(4).
\emph{IEEE Trans.\ Inform.\ Theory}, 44(4):1369--1387, 1998.

\bibitem[DKS17]{DKS17}
I.~Diakonikolas, D.~M.~Kane, and A.~Stewart.
Statistical query lower bounds for robust estimation of high-dimensional Gaussians and Gaussian mixtures.
In \emph{FOCS}, pages 73--84. IEEE, 2017.

\bibitem[DGM14]{DGM14}
N.~D\"ottling, J.~M\"uller-Quade, and A.~Nascimento.
IND-CCA secure cryptography based on a variant of the LPN problem.
In \emph{ASIACRYPT}, pages 485--508. Springer, 2014.

\bibitem[DP12]{DP12}
Y.~Dodis, K.~Pietrzak, and D.~Wichs.
Key derivation without entropy waste.
In \emph{EUROCRYPT}, pages 93--110. Springer, 2012.

\bibitem[EKM17]{EKM17}
A.~Esser, R.~K\"ubler, and A.~May.
LPN decoded.
In \emph{CRYPTO}, pages 486--514. Springer, 2017.

\bibitem[FO99]{FO99}
E.~Fujisaki and T.~Okamoto.
Secure integration of asymmetric and symmetric encryption schemes.
In \emph{CRYPTO}, pages 537--554. Springer, 1999.

\bibitem[FGR+17]{FGR+17}
V.~Feldman, E.~Grigorescu, L.~Reyzin, and S.~Vempala.
Statistical query lower bounds for robust estimation of high-dimensional Gaussians and Gaussian mixtures.
In \emph{FOCS}, pages 211--220. IEEE, 2017.

\bibitem[FV19]{FV19}
A.~Fazeli and A.~Vardy.
On the scaling exponent of binary polarization kernels.
In \emph{IEEE ISIT}, pages 1472--1476, 2019.

\bibitem[GGPR13]{GGPR13}
R.~Gennaro, C.~Gentry, B.~Parno, and M.~Raykova.
Quadratic span programs and succinct NIZKs without PCPs.
In \emph{EUROCRYPT}, pages 626--645. Springer, 2013.

\bibitem[GJL14]{GJL14}
Q.~Guo, T.~Johansson, and C.~L\"ondahl.
Solving LPN using covering codes.
In \emph{ASIACRYPT}, pages 1--20. Springer, 2014.

\bibitem[GL22]{GL22}
S.~Gollakota and D.~Liang.
PAC-learning noisy stabilizer states in the statistical query model.
\emph{arXiv preprint}, 2022.

\bibitem[Got97]{Got97}
D.~Gottesman.
Stabilizer codes and quantum error correction.
PhD thesis, Caltech, 1997.

\bibitem[Gro96]{Gro96}
L.~K.~Grover.
A fast quantum mechanical algorithm for database search.
In \emph{STOC}, pages 212--219. ACM, 1996.

\bibitem[Gro16]{Gro16}
J.~Groth.
On the size of pairing-based non-interactive arguments.
In \emph{EUROCRYPT}, pages 305--326. Springer, 2016.

\bibitem[HHK17]{HHK17}
D.~Hofheinz, K.~H\"ovelmanns, and E.~Kiltz.
A modular analysis of the Fujisaki--Okamoto transformation.
In \emph{TCC}, pages 341--371. Springer, 2017.

\bibitem[HKL+12]{HKL+12}
S.~Heyse, E.~Kiltz, V.~Lyubashevsky, C.~Paar, and K.~Pietrzak.
Lapin: An efficient authentication protocol based on ring-LPN.
In \emph{FSE}, pages 346--365. Springer, 2012.

\bibitem[HVU14]{HVU14}
S.~H.~Hassani, R.~Urbanke, and S.~V.~Macris.
Near-optimal finite-length scaling for polar codes over large alphabets.
\emph{IEEE Trans.\ Inform.\ Theory}, 60(11):6960--6976, 2014.

\bibitem[Kea98]{Kea98}
M.~Kearns.
Efficient noise-tolerant learning from statistical queries.
\emph{J.\ ACM}, 45(6):983--1006, 1998.

\bibitem[Kir11]{Kir11}
P.~Kirchner.
Improved generalized birthday attack.
\emph{Cryptology ePrint Archive}, Report 2011/377, 2011.

\bibitem[KLP+25]{KLP+25}
A.~Khesin, J.~Lu, A.~Poremba, S.~Ramkumar, and V.~Vaikuntanathan.
Hardness of decoding random quantum stabilizer codes.
\emph{arXiv:2509.20697}, 2025.

\bibitem[LF06]{LF06}
E.~Levieil and P.-A.~Fouque.
An improved LPN algorithm.
In \emph{SCN}, pages 348--359. Springer, 2006.

\bibitem[LPQR26]{LPQR26}
J.~Lu, A.~Poremba, Y.~Quek, and S.~Ramkumar.
Post-quantum cryptography from quantum stabilizer decoding.
\emph{arXiv:2603.19110}, 2026.

\bibitem[Lyu05]{Lyu05}
V.~Lyubashevsky.
The parity problem in the presence of noise, decoding random linear codes, and the subset sum problem.
In \emph{APPROX-RANDOM}, pages 378--389. Springer, 2005.

\bibitem[Mac71]{Mac71}
I.~G.~Macdonald.
Symmetric Functions and Hall Polynomials.
Oxford University Press, 1971.

\bibitem[McE78]{McE78}
R.~J.~McEliece.
A public-key cryptosystem based on algebraic coding theory.
\emph{DSN Progress Report}, 44:114--116, 1978.

\bibitem[MT09]{MT09}
R.~Mori and T.~Tanaka.
Performance and construction of polar codes on symmetric binary-input memoryless channels.
In \emph{IEEE ISIT}, pages 1496--1500, 2009.

\bibitem[NC00]{NC00}
M.~A.~Nielsen and I.~L.~Chuang.
Quantum Computation and Quantum Information.
Cambridge University Press, 2000.

\bibitem[NIS22]{NIS22}
NIST.
Post-quantum cryptography standardization, 2022.
\url{https://csrc.nist.gov/projects/post-quantum-cryptography}.

\bibitem[PHGR13]{PHGR13}
B.~Parno, J.~Howell, C.~Gentry, and M.~Raykova.
Pinocchio: Nearly practical verifiable computation.
In \emph{IEEE S\&P}, pages 238--252, 2013.

\bibitem[PQS26]{PQS26}
A.~Poremba, Y.~Quek, and P.~Shor.
Learning stabilizers with noise.
\emph{arXiv:2410.18953}, 2025.

\bibitem[Pra62]{Pra62}
E.~Prange.
The use of information sets in decoding cyclic codes.
\emph{IRE Trans.\ Inform.\ Theory}, 8(5):5--9, 1962.

\bibitem[Reg05]{Reg05}
O.~Regev.
On lattices, learning with errors, random linear codes, and cryptography.
\emph{J.\ ACM}, 56(6):34, 2005.

\bibitem[SAB+22]{SAB+22}
P.~Schwabe, R.~Avanzi, J.~Bos, L.~Ducas, E.~Kiltz, T.~Lepoint, V.~Lyubashevsky, J.~M.~Schanck, G.~Seiler, and D.~Stehl\'e.
CRYSTALS-Kyber.
NIST PQC Round 3 Submission, 2022.

\bibitem[Sei15]{Sei15}
M.~Seidl.
Polar coding: Finite-length aspects.
PhD thesis, Friedrich-Alexander-Universit\"at Erlangen-N\"urnberg, 2015.

\bibitem[Sho94]{Sho94}
P.~W.~Shor.
Algorithms for quantum computation: Discrete logarithms and factoring.
In \emph{FOCS}, pages 124--134. IEEE, 1994.

\bibitem[Tri12]{Tri12}
P.~Trifonov.
Efficient design and decoding of polar codes.
\emph{IEEE Trans.\ Commun.}, 60(11):3221--3227, 2012.

\bibitem[TV15]{TV15}
I.~Tal and A.~Vardy.
List decoding of polar codes.
\emph{IEEE Trans.\ Inform.\ Theory}, 61(5):2213--2226, 2015.

\bibitem[WS21]{WS21}
A.~Wang and F.~Zhang.
Revisiting the BKW algorithm for LPN.
\emph{Cryptology ePrint Archive}, Report 2021/213, 2021.

\bibitem[ZJW16]{ZJW16}
B.-Y.~Zhang, L.~Jia, and M.-W.~Wang.
Security analysis of LPN under perfect code cover.
\emph{IEEE Access}, 4:3979--3986, 2016.

\bibitem[komm]{komm}
A.~Krohn.
\texttt{komm}: An open-source library for communication systems.
\url{https://github.com/ryukau/komm}, 2024.

\end{thebibliography}

\end{document}
